\chapter{Dupuytren Contracture Release}

\section*{Introduction \& Clinical Indications}
Dupuytren contracture release is a surgical intervention designed to correct the progressive, often painless, flexion contracture of the fingers and thumb caused by Dupuytren's disease. This condition involves the pathological thickening and shortening of the palmar fascia, forming palpable nodules and cords that restrict the normal extension of the digits, most commonly the ring and small fingers. The primary goal of the surgery is to excise or divide these fibrotic cords, thereby releasing the contracted metacarpophalangeal (MCP) and proximal interphalangeal (PIP) joints and restoring functional finger extension. This procedure is clinically significant as it can dramatically improve hand function, enabling patients to perform daily activities that were previously hindered by the contracture and enhancing their overall quality of life.

\begin{itemize}
  \item Dupuytren's Fasciectomy
  \item Selective Fasciectomy
  \item Regional Fasciectomy
  \item Palmar Fascia Excision
  \item Open Fasciotomy
  \item Aponeurectomy
\end{itemize}

The fundamental surgical principle of Dupuytren contracture release is the mechanical interruption or removal of the pathological collagenous cords within the palmar aponeurosis that are responsible for the fixed flexion deformity of the digits. These cords, composed primarily of type III collagen and myofibroblasts, shorten and thicken over time, tethering the skin and underlying structures and preventing full finger extension. The primary technical goal is to release this tethering effect, thereby restoring the normal length of the fascial structures and allowing the affected joints to extend freely.

The rationale for surgical intervention stems from the progressive nature of the disease and its significant impact on hand function. By excising (fasciectomy) or dividing (fasciotomy) the diseased tissue, the surgeon aims to alleviate the mechanical obstruction to finger extension. A fasciectomy involves the meticulous excision of the diseased cords and associated nodules, aiming for a more complete removal of the pathological tissue, which is thought to reduce the risk of recurrence in the immediate area. Conversely, a fasciotomy involves the simple division of the contracted cords, either openly or percutaneously, without excising tissue. Both techniques aim to improve the functional arc of motion of the hand, with careful dissection paramount to avoid injury to the adjacent neurovascular bundles, which can be intimately involved with or displaced by the contracting cords.

The condition now known as Dupuytren's disease was first comprehensively described by Baron Guillaume Dupuytren, a distinguished French surgeon, in 1831. He presented his findings at the Royal Academy of Medicine in Paris and performed the first recorded fasciotomy, dividing the contracted cords through small incisions. His initial technique involved a longitudinal incision and division of the cords, which was a a significant advancement at the time, moving beyond amputation as the only option.

Throughout the late 19th and early 20th centuries, surgical approaches evolved. Surgeons like Sir Astley Cooper and others contributed to the understanding and treatment of the disease, refining Dupuytren's original technique. The concept of fasciectomy, involving the excision of the diseased tissue, gained prominence in the mid-20th century, particularly with the work of Skoog, who advocated for a more radical removal of the diseased fascia to reduce recurrence. This led to the development of regional and selective fasciectomy techniques. More recently, less invasive techniques have emerged, such as percutaneous needle aponeurotomy (PNA), popularized in the 1990s, which offers a minimally invasive alternative by using a needle to divide the cords. The introduction of collagenase clostridium histolyticum injections in the early 21st century provided a non-surgical option, enzymatically dissolving the collagen cords. These advancements reflect a continuous effort to balance efficacy, invasiveness, and patient recovery in the management of Dupuytren's contracture.

Surgical intervention for Dupuytren's contracture is typically indicated when the disease significantly impairs hand function or is rapidly progressing.

\begin{itemize}
  \item When a Dupuytren cord causes a functional contracture, typically a metacarpophalangeal (MCP) joint contracture of 30 degrees or more, as measured by goniometry.
  \item For any degree of proximal interphalangeal (PIP) joint contracture, as these are often more functionally limiting and challenging to correct non-surgically.
  \item When the contracture results in a positive "tabletop test," meaning the patient cannot lay their palm flat on a surface.
  \item When the contracture significantly interferes with daily activities such as washing, dressing, shaking hands, grasping objects, or occupational tasks.
  \item In cases of rapid progression of the contracture, even if the absolute degree of contracture is not yet severe, to prevent further functional loss.
  \item When pain is associated with the nodules or cords, although this is a less common primary indication for surgery.
  \item For recurrent contractures following previous non-surgical or surgical interventions, if functional impairment persists and warrants further treatment.
  \item When the contracture affects the thumb or multiple digits, leading to complex functional deficits that require comprehensive correction.
\end{itemize}

Dupuytren contracture release is considered a primary treatment for disabling Dupuytren's contracture that significantly impacts hand function. It is typically offered when non-surgical options, such as observation for mild, non-progressive cases or collagenase injections for specific cord types, are either not indicated, have failed to achieve adequate correction, or are less likely to provide a durable solution. For patients presenting with established, functionally limiting contractures, especially those involving the proximal interphalangeal (PIP) joints or severe metacarpophalangeal (MCP) joint contractures, surgical release is often the first-line intervention due to its efficacy in achieving significant and lasting correction. In cases of recurrence after initial treatment, whether surgical or non-surgical, a repeat surgical release may be considered a secondary or salvage procedure, particularly if the contracture is severe or rapidly progressive. The choice between primary surgical intervention and other modalities depends on the severity and location of the contracture, the patient's overall health, their functional demands, and the surgeon's expertise.

\section*{Relevant Surgical Anatomy}
The hand's intricate anatomy is crucial for understanding Dupuytren's contracture and its surgical release. The palmar aponeurosis, a dense fibrous sheet located beneath the skin of the palm, is the primary structure affected. This aponeurosis extends distally to form pretendinous bands, which then divide into spiral bands, lateral bands, and natatory ligaments, ultimately inserting into the dermis and phalanges. In Dupuytren's disease, these fascial structures undergo fibrotic transformation, forming nodules and cords that shorten and thicken, tethering the skin and underlying tendons, thus preventing full digital extension.

Crucially, the digital nerves and arteries run in close proximity to, and can become intimately involved with or displaced by, these contracting cords. The common digital nerves and arteries bifurcate into proper digital nerves and arteries, which course along the sides of the fingers. These neurovascular bundles are highly vulnerable to iatrogenic injury during dissection. The flexor tendons lie deep to the palmar fascia, and while not directly affected by the disease process, their function is mechanically impeded by the overlying contracture. Surgical landmarks include the distal palmar crease, proximal digital creases, and the precise location of the palpable cords and nodules, which guide incision planning and dissection. Lymphatic drainage of the hand primarily follows the venous system, ascending to the epitrochlear and axillary lymph nodes.

\section*{Preoperative Workup \& Preparation}
Preoperative preparation for Dupuytren contracture release involves a thorough clinical assessment and comprehensive patient education to ensure optimal outcomes and minimize risks. A detailed hand examination is performed to accurately measure the degree of contracture in each affected joint using a goniometer, assess neurovascular status, and document skin quality and the presence of nodules or pits. Preoperative photographs are often taken to document the baseline condition.

\begin{itemize}
  \item \textbf{Medical Evaluation:} A general medical evaluation, including cardiac and pulmonary assessment, is conducted to ensure the patient is fit for anesthesia and surgery.
  \item \textbf{Laboratory Tests:} Routine preoperative laboratory tests, such as a complete blood count, basic metabolic panel, and coagulation profile, are typically ordered, especially if regional or general anesthesia is planned.
  \item \textbf{Medication Review:} Patients are advised to discontinue anticoagulants or antiplatelet medications (e.g., aspirin, warfarin, novel oral anticoagulants) several days prior to surgery, as directed by their surgeon and primary care physician, to minimize bleeding risk.
  \item \textbf{NPO Status:} Patients are instructed to be NPO (nothing by mouth) for a specified period (typically 6-8 hours) before surgery, depending on the type of anesthesia.
  \item \textbf{Prophylactic Antibiotics:} Intravenous prophylactic antibiotics are usually administered just before the incision to reduce the risk of surgical site infection.
  \item \textbf{Informed Consent:} A comprehensive informed consent process is crucial, covering the risks, benefits, alternatives (including non-surgical options), the potential for disease recurrence, and the importance of postoperative hand therapy.
\end{itemize}

Careful consideration of contraindications and precautions is essential to ensure patient safety and optimize surgical outcomes for Dupuytren contracture release.

\textbf{Situations requiring treatment, optimization, or an alternative plan:}
\begin{itemize}
  \item Active infection in the hand or surrounding tissues, which must be treated prior to surgery to prevent spread.
  \item Severe, uncontrolled systemic comorbidities (e.g., unstable angina, recent myocardial infarction, uncontrolled diabetes, severe coagulopathy) that significantly increase surgical and anesthetic risk.
  \item Patient refusal or inability to comprehend and comply with the demanding postoperative rehabilitation protocols.
  \item Mild contractures that do not cause functional impairment, are not progressing rapidly, and do not meet established criteria for intervention.
\end{itemize}

\textbf{Relative Contraindications \& Precautions:}
\begin{itemize}
  \item Very tight or recurrent contractures, especially those involving the proximal interphalangeal (PIP) joints, which carry a higher risk of neurovascular injury due to distorted and scarred anatomy.
  \item Poor skin quality, such as thin, atrophic skin, active ulceration, or previous radiation therapy, which may compromise wound healing and increase the risk of skin necrosis.
  \item Severe peripheral vascular disease, which can impair wound healing, increase infection risk, and compromise digital viability.
  \item Uncontrolled diabetes mellitus, which can increase infection risk, impair wound healing, and contribute to microvascular complications.
  \item Previous failed surgery with significant scarring, making subsequent dissection more challenging and increasing the risk of iatrogenic injury.
  \item Patients with a history of Complex Regional Pain Syndrome (CRPS), as surgery can sometimes trigger or exacerbate this debilitating condition.
  \item The digital nerves and arteries may be intimately involved with or displaced by the disease cords, requiring meticulous and careful dissection to prevent iatrogenic injury, often necessitating the use of magnification.
  \item Active smoking, which significantly impairs wound healing and increases the risk of complications; smoking cessation is strongly advised preoperatively.
\end{itemize}

\section*{Surgical Technique \& Procedure}
The procedure is typically performed under regional anesthesia (e.g., axillary block) or local anesthesia with sedation, often combined with a tourniquet applied to the upper arm to ensure a bloodless field. The patient is positioned supine with the arm extended on an arm board.

\begin{itemize}
  \item \textbf{Anesthesia \& Positioning:} Regional block (axillary or supraclavicular) or general anesthesia is administered. The patient is supine, and the affected arm is placed on an arm board. A tourniquet is inflated on the upper arm to create a bloodless field, optimizing visualization.
  \item \textbf{Incision Planning:} Zigzag (Bruner) or transverse incisions are meticulously planned over the affected palmar skin and digits. These designs are chosen to maximize exposure while minimizing the risk of postoperative scar contracture.
  \item \textbf{Skin Flap Elevation:} The skin flaps are carefully elevated, exposing the underlying diseased palmar aponeurosis. This step requires precision to avoid injury to the superficial neurovascular structures.
  \item \textbf{Cord Identification \& Dissection:} The fibrotic Dupuytren's cords are meticulously identified. Using loupe magnification, the surgeon carefully dissects these cords free from the surrounding healthy tissue, paying particular attention to the delicate digital nerves and arteries that may be encased within or displaced by the cords.
  \item \textbf{Cord Excision/Division:} For a fasciectomy, the diseased cords and associated nodules are excised, either selectively (removing only the visibly diseased tissue) or regionally (removing a broader area of fascia). For a fasciotomy, the cords are simply divided using a scalpel or needle, without extensive tissue removal.
  \item \textbf{Hemostasis \& Tourniquet Release:} Once the cords are addressed, meticulous hemostasis is achieved using electrocautery. The tourniquet is then deflated, and any remaining bleeding points are controlled.
  \item \textbf{Wound Closure \& Dressing:} The wound is irrigated, and the skin incisions are closed with fine sutures. In cases of significant skin deficit or severe contracture, a Z-plasty or skin graft may be necessary to achieve full extension and prevent skin contracture. A bulky dressing and a custom splint are applied to maintain the corrected finger position and protect the surgical site.
\end{itemize}

\section*{Complications \& Risks}
While Dupuytren contracture release is generally safe and effective, like any surgical procedure, it carries potential risks and complications that patients should be aware of.

\begin{itemize}
  \item \textbf{General Surgical Complications:}
  \begin{itemize}
    \item \textit{Bleeding and Hematoma Formation:} Can cause pain, swelling, and pressure on nerves, potentially requiring evacuation.
    \item \textit{Infection of the Surgical Site:} Ranging from superficial cellulitis to deep space infection, potentially requiring antibiotics or further surgical debridement.
    \item \textit{Anesthesia Complications:} Including allergic reactions, respiratory issues, cardiovascular events (e.g., myocardial infarction, stroke), or nerve injury related to regional blocks.
    \item \textit{Scarring:} Hypertrophic or keloid scarring, which can be cosmetically unappealing and cause discomfort or functional limitation.
  \end{itemize}
  \item \textbf{Procedure-Specific Risks:}
  \begin{itemize}
    \item \textit{Nerve Injury:} Leading to temporary or permanent numbness, tingling, or weakness in the affected digit. The proper digital nerves are particularly vulnerable due to their close proximity to the cords.
    \item \textit{Vascular Injury:} Potentially compromising blood supply to the finger, which can lead to tissue necrosis, delayed wound healing, or even digit loss in severe cases.
    \item \textit{Skin Necrosis or Wound Healing Problems:} Especially with extensive skin dissection, in patients with poor circulation, or those with compromised skin quality.
    \item \textit{Incomplete Correction of the Contracture:} Particularly challenging in severe or long-standing proximal interphalangeal (PIP) joint contractures, where full extension may not be achievable.
    \item \textit{Joint Stiffness or Limited Range of Motion:} Despite successful cord release, often requiring intensive and prolonged hand therapy to regain optimal mobility.
    \item \textit{Disease Recurrence:} As Dupuytren's is a progressive condition, the underlying diathesis remains. Recurrence rates vary by technique, disease severity, and patient factors, necessitating long-term follow-up.
    \item \textit{Complex Regional Pain Syndrome (CRPS):} A rare but severe chronic pain condition that can affect the hand and arm, characterized by disproportionate pain, swelling, and trophic changes.
    \item \textit{Tendon Injury:} Though rare, can occur during meticulous dissection of deeply adherent cords, potentially requiring repair.
    \item \textit{Digital Nerve Neuroma:} A painful benign growth on a nerve, which can result from nerve irritation or injury during surgery.
    \item \textit{Flap Necrosis or Graft Failure:} If skin grafts or local flaps are used to cover skin defects, these may fail to take or necrose.
  \end{itemize}
\end{itemize}

\section*{Postoperative Care \& Recovery}
Immediately after Dupuytren contracture release, the hand is typically placed in a bulky dressing and a splint to maintain the corrected finger extension and protect the surgical site. The patient will spend time in the post-anesthesia care unit (PACU) for monitoring of vital signs, pain management, and neurovascular checks of the digits. Elevation of the hand above heart level is crucial in the initial 24-48 hours to minimize swelling and reduce pain. Oral pain medication will be prescribed to manage postoperative discomfort, and patients are encouraged to begin gentle active range of motion exercises of the unaffected digits.

Over the first 1-2 weeks, the patient will begin a structured hand therapy program, which is vital for optimizing outcomes. This typically involves gentle active and passive range of motion exercises to prevent stiffness, scar massage, and edema control techniques. Sutures are usually removed around 10-14 days post-surgery. Night splinting, often custom-made to maintain extension, may be recommended for several months to prevent contracture recurrence during sleep. Patients are gradually encouraged to resume light daily activities, with a progressive increase in activity level under the guidance of a hand therapist. Full recovery of strength and function can take several months, often 3-6 months, and adherence to the rehabilitation program is critical for achieving the best possible long-term results and minimizing the risk of stiffness or recurrence.

During Dupuytren contracture release, the surgeon's findings typically include the presence of thickened, fibrous cords and nodules within the palmar aponeurosis, often extending into the digits. These cords are usually firm, pale, and may be intimately associated with or displacing the digital neurovascular bundles, requiring careful dissection. The degree of contracture in the metacarpophalangeal (MCP) and proximal interphalangeal (PIP) joints is visually assessed and documented intraoperatively, along with the extent of skin tethering or pitting. Any intraoperative complications, such as nerve or vessel injury, are also noted.

If a fasciectomy is performed, the excised tissue is sent for histopathological examination. The pathology report will confirm the diagnosis of Dupuytren's disease, characterized by hypercellular fibrous tissue composed predominantly of myofibroblasts and disorganized type III collagen. These findings are consistent with the proliferative and involutional phases of the disease, validating the surgical indication and confirming the nature of the resected tissue. The presence of inflammatory cells is usually minimal, distinguishing it from other fibrotic conditions.

A normal and successful postoperative course following Dupuytren contracture release is characterized by a significant and sustained improvement in the ability to extend the affected finger(s), often achieving a straight or near-straight digit. This correction should translate into restored hand function, allowing the patient to comfortably perform tasks that were previously difficult or impossible, such as placing their hand flat on a table, washing their face, or shaking hands. Postoperative wound healing should be uncomplicated, with minimal scarring and no signs of infection or skin necrosis. The patient should experience a gradual reduction in pain and swelling, which typically subsides within the first few weeks. With diligent adherence to the prescribed hand therapy program, patients are expected to regain a good range of motion and strength in the hand, progressively returning to normal daily activities and eventually more strenuous tasks over several months. The ultimate goal is a functional, pain-free hand with improved dexterity and quality of life.

Postoperative care for Dupuytren contracture release is critical for optimizing outcomes and preventing complications. It involves a multidisciplinary approach focused on wound healing, pain management, and functional rehabilitation.

\begin{itemize}
  \item \textbf{Postoperative Clinic Visits:} Regular appointments with the surgeon are scheduled, typically at 1-2 weeks, 6 weeks, 3 months, and 6-12 months, to monitor wound healing, assess range of motion, evaluate overall recovery, and screen for early signs of recurrence.
  \item \textbf{Suture Removal:} Sutures are typically removed 10-14 days after surgery, often during the first follow-up visit.
  \item \textbf{Hand Therapy/Physical Therapy:} This is an essential component of recovery, commencing within days of surgery. It includes active and passive range of motion exercises, scar management techniques (e.g., massage, silicone sheeting), edema control, and progressive strengthening exercises. Custom splinting, especially night splints to maintain extension, is often prescribed for several months to prevent contracture recurrence.
  \item \textbf{Activity Resumption Guidance:} Patients receive specific instructions on a gradual return to normal daily activities and work, with initial restrictions on heavy gripping, lifting, or strenuous activities to protect the healing tissues.
  \item \textbf{Surveillance for Recurrence:} Ongoing monitoring for signs of new nodules or cords, or the return of contracture, is crucial, as Dupuytren's disease is progressive. Patients are educated on self-monitoring and when to seek further medical advice.
\end{itemize}
Patients are advised to contact their surgeon immediately if they experience increased pain, redness, swelling, fever, persistent numbness or tingling, foul-smelling discharge from the wound, or an inability to move their fingers as instructed.

\section*{Outcomes \& Prognosis}
Dupuytren's disease, the underlying condition necessitating contracture release, exhibits a strong genetic predisposition and is most prevalent in individuals of Northern European (often referred to as "Viking") descent. Its incidence significantly increases with age, typically affecting individuals over 50 years old, and is notably more common in males than females, with a male-to-female ratio often cited as approximately 7:1. The ring and small fingers are most frequently affected, followed by the middle finger, and less commonly the thumb or index finger. While the exact prevalence varies geographically, studies suggest rates as high as 20-30\% in some populations of Northern European ancestry. Associated risk factors include diabetes mellitus, epilepsy, heavy alcohol consumption, smoking, and manual labor, though the exact causal relationships are complex. Although Dupuytren's disease itself is not life-threatening, the morbidity associated with severe contractures can significantly impair hand function and quality of life. Surgical intervention carries a very low mortality rate, but morbidity can arise from complications such as nerve injury, infection, and recurrence, impacting functional outcomes.

The outcomes of Dupuytren contracture release are generally favorable, with high initial success rates in correcting finger contractures and significantly improving hand function. For metacarpophalangeal (MCP) joint contractures, surgical release typically achieves excellent correction, often restoring full extension. Proximal interphalangeal (PIP) joint contractures are inherently more challenging to correct, and while significant improvement is common, complete correction may not always be achieved, especially in long-standing or severe cases with secondary joint changes. Functional outcomes, as measured by patient-reported questionnaires (e.g., QuickDASH) and objective hand assessments, consistently demonstrate substantial improvement in the ability to perform daily activities and overall quality of life.

Dupuytren's disease can recur and may progress in other cords or digits. Recurrence varies with disease biology, joint, severity, technique, and follow-up; less invasive procedures may recover faster but can recur more often than fasciectomy. Patients should be counseled that PIP contractures and long-standing disease are harder to fully correct.

\section*{References}
\begin{enumerate}
  \item Schwartz SI, Brunicardi FC, Andersen DK, et al. Schwartz's Principles of Surgery. 11th ed. McGraw-Hill Education; 2019.
  \item Green DP, Hotchkiss RN, Pederson WC, Wolfe SW. Green's Operative Hand Surgery. 7th ed. Elsevier; 2017.
  \item American Society for Surgery of the Hand. Dupuytren's Contracture. Available at: \texttt{https://www.assh.org/handcare/condition/dupuytrens-contracture}. Accessed 2024.
  \item MedlinePlus. Dupuytren contracture surgery. U.S. National Library of Medicine; 2025.
\end{enumerate}

\clearpage
